% Part: lambda-calculus
% Chapter: introduction
% Section: beta

\documentclass[../../../include/open-logic-section]{subfiles}

\begin{document}

\olfileid[id]{lam}{int}{bet}
\olsection{Reduksi-$\beta$}

Ketika melihat $(\lambd[m][(\lambd[y][y]) m])$, wajar untuk menduga bahwa
suku itu mempunyai suatu hubungan dengan $\lambd[m][m]$, yakni bahwa suku
kedua semestinya merupakan hasil ``penyederhanaan'' suku pertama. Gagasan
\emph{reduksi-$\beta$} menangkap intuisi ini secara formal.

\begin{defn}[Kontraksi-$\beta$, $\bredone$] \ollabel{defn:betacontr}
  \emph{Kontraksi-$\beta$} ($\bredone$) adalah relasi kompatibel terkecil pada
  suku-suku yang memenuhi syarat berikut:
  \[
    (\lambd[x][N])Q \bredone \Subst{N}{Q}{x}
  \]
  Kita mengatakan bahwa $P$ \emph{dikontraksi-$\beta$} menjadi $Q$ jika
  $P \bredone Q$. Suku berbentuk $(\lambd[x][N])Q$ disebut \emph{redex}.
\end{defn}

\begin{prob} \ollabel{prob:def}
  Tuliskan secara lengkap definisi induktif ekuivalen dari kontraksi-$\beta$,
  sebagaimana yang kita lakukan untuk pengubahan variabel terikat dalam
  \olref[lam][syn][alp]{defn:aconvone}.
\end{prob}

\begin{defn}[Reduksi-$\beta$, $\bred$] \ollabel{defn:betared}
  \emph{Reduksi-$\beta$} ($\bred$) adalah relasi refleksif dan transitif
  terkecil pada suku-suku yang memuat $\bredone$. Kita mengatakan bahwa $P$
  tereduksi-$\beta$ menjadi $Q$ jika $P \bred Q$.
\end{defn}

Kita akan menulis $\redone$ sebagai ganti $\bredone$, dan $\red$ sebagai ganti
$\bred$, apabila konteksnya jelas.

Secara informal, $M \bred N$ jika dan hanya jika $M$ dapat diubah menjadi $N$
melalui nol atau beberapa langkah kontraksi-$\beta$.

\begin{defn}[Normal-$\beta$]
Suku yang tidak dapat dikontraksi-$\beta$ lebih lanjut disebut
\emph{normal-$\beta$}.
\end{defn}

Jika $M \bred N$ dan $N$ normal-$\beta$, kita mengatakan bahwa $N$ merupakan
\emph{bentuk normal} dari~$M$. Kita dapat bertanya apakah bentuk normal suatu
suku unik; jawabannya adalah ya, sebagaimana akan kita lihat nanti.

Mari kita perhatikan beberapa contoh.
\begin{enumerate}
\item Kita mempunyai
\begin{align*}
(\lambd[x][xxy]) \lambd[z][z] & \redone (\lambd[z][z])(\lambd[z][z]) y \\
& \redone (\lambd[z][z]) y \\
& \redone y
\end{align*}
\item ``Menyederhanakan'' suatu suku justru dapat membuatnya lebih rumit:
\begin{align*}
(\lambd[x][xxy])(\lambd[x][xxy]) & \redone (\lambd[x][xxy])(\lambd[x][xxy])y \\
& \redone (\lambd[x][xxy])(\lambd[x][xxy])yy \\
& \redone \dots
\end{align*}
\item Penyederhanaan juga dapat membiarkan suatu suku tidak berubah:
\[
(\lambd[x][xx])(\lambd[x][xx]) \redone (\lambd[x][xx])(\lambd[x][xx])
\]
\item Selain itu, beberapa suku dapat direduksi dengan lebih dari satu cara;
  sebagai contoh,
\[
(\lambd[x][(\lambd[y][yx]) z)] v \redone (\lambd[y][yv]) z
\]
dengan mengontraksi aplikasi terluar; dan
\[
(\lambd[x][(\lambd[y][yx]) z)] v \redone (\lambd[x][zx]) v
\]
dengan mengontraksi aplikasi terdalam. Namun, perhatikan bahwa dalam kasus ini
kedua suku selanjutnya tereduksi menjadi suku yang sama, yakni $zv$.
\end{enumerate}

Hasil akhir dalam contoh terakhir bukanlah suatu kebetulan, melainkan
memperlihatkan sifat kalkulus lambda yang mendalam dan penting, yang dikenal
sebagai sifat Church--Rosser.

\begin{digress}
  Secara umum, terdapat lebih dari satu cara untuk mereduksi-$\beta$ suatu
  suku. Karena itu, banyak strategi reduksi telah ditemukan, dan yang paling
  lazim di antaranya adalah \emph{strategi alami}. Strategi alami selalu
  mengontraksi redex \emph{paling kiri}, dengan posisi redex didefinisikan
  sebagai titik awalnya dalam suku. Strategi alami mempunyai sifat berguna
  bahwa suatu suku dapat direduksi ke bentuk normal dengan suatu strategi jika
  dan hanya jika suku itu dapat direduksi ke bentuk normal menggunakan
  strategi alami. Selanjutnya kita akan menggunakan strategi alami kecuali
  jika dinyatakan lain.
\end{digress}

\begin{defn}[Ekuivalensi-$\beta$, $\equal$]
  \emph{Ekuivalensi-$\beta$} ($\equal$) adalah relasi yang didefinisikan
  secara induktif sebagai berikut:
  \begin{enumerate}
  \item $M \equal M$.
  \item Jika $M \equal N$, maka $N \equal M$.
  \item Jika $M \equal N$, $N \equal O$, maka $M \equal O$.
  \item Jika $M \equal N$, maka $PM \equal PN$.
  \item Jika $M \equal N$, maka $MQ \equal NQ$.
  \item Jika $M \equal N$, maka $\lambd[x][M] \equal \lambd[x][N]$.
  \item $(\lambd[x][N])Q \equal \Subst{N}{Q}{x}$.
  \end{enumerate}
\end{defn}

Tiga aturan pertama menjadikan relasi tersebut suatu relasi ekuivalensi; tiga
aturan berikutnya membuatnya kompatibel; aturan terakhir memastikan bahwa
relasi tersebut memuat kontraksi-$\beta$.

Secara informal, dua suku $\beta$-ekuivalen jika dan hanya jika salah satunya
dapat diubah menjadi yang lain melalui nol atau beberapa langkah
kontraksi-$\beta$, atau ``kebalikan'' kontraksi-$\beta$. Kebalikan
kontraksi-$\beta$ didefinisikan sedemikian sehingga $M$ dikontraksi-$\beta$
secara terbalik menjadi $N$ jika dan hanya jika $N$ dikontraksi-$\beta$
menjadi $M$.

Selain aturan-aturan di atas, kita akan memperluas relasi tersebut dengan
aturan-aturan lain, dan menyatakan relasi ekuivalensi yang diperluas itu
sebagai $\equal[X]$, dengan $X$ aturan yang memperluasnya.

\end{document}
